\section{Introduction}
\label{sec:intro}

The cost of retrieval failure is clearest when the requested source is
missing. In our demonstration, a user asks which training epochs appear
in a fictitious paper. The collection returns generic BERT recipes and
weak title matches, but no evidence about the named work. Figure~\ref{fig:teaser}
shows the platform reporting what it searched and abstaining instead of
guessing an epoch. This case motivates a broader systems problem:
knowledge-base quality depends on where documents are placed, whether
parsing and indexing succeeded, and whether an answer can be traced to
the passage that supports it.

Retrieval-augmented generation (RAG) commonly ranks passages and places
them in a generation context~\cite{rag-lewis}. Self-RAG adds model-side
reflection~\cite{selfrag}, while Corrective RAG evaluates retrieved
documents~\cite{crag}. These methods address evidence selection and
generation; they do not, by themselves, specify how a shared collection
is organized and kept consistent with the addresses later used for
inspection. Hierarchical and graph-based systems organize retrieval
representations in different ways~\cite{raptor,deepread,graphrag,lightrag}.
Our contribution is a system integration and interactive demo, not a new
retrieval model: collection operations, retrieval, and citations use the
same durable address space.

The platform makes four design choices. \emph{One address space}: each
source is represented as base, document, part, and section throughout
storage, retrieval, and citation. \emph{Gated organization}: ingestion
parses, deduplicates, splits, routes, describes, tags, and indexes
documents, then checks content, metadata, and index consistency.
\emph{Separate recall from evidence judgment}: a layered search proposes
candidates, and a local engine decides whether each segment contains
answerable evidence. \emph{Report failure explicitly}: when no candidate
survives, the agent returns a scoped not-found report with the search
path and blind spots.

Keeping the file-backed hierarchy, metadata, vector index, optional
graph, and agent-facing tools consistent across ingest and retrieval is
the main engineering challenge. The demo exposes these checks to
research administrators and agent developers, an audience building
shared information services for CIKM's data-management and retrieval
community. Sections~\ref{sec:system}--\ref{sec:eval} describe the
platform, attendee workflow, and measured limits of the retrieval
mechanism.
